\documentclass[../main.tex]{subfiles}

\begin{document}
\subsection{Lema de Goursat}
\classdate{2026-04-14}

\begin{lemma}[Lema de Goursat]\label{lemma:lem-goursat-v1}
    Sean \(f \colon \Omega \subseteq \mathbb{C} \to \mathbb{C}\) una función analítica
    (en \(\Omega\)) y \(R = [a, b] \mul [c, d] \subseteq \Omega\), entonces

    \begin{equation*}
        \int_{\gamma}^{}f(z)\odif{z} = 0
    \end{equation*}

    donde \(\gamma\) es una parametrización de \(\partial R =
    \mathrm{Fr}(R)\) que la recorre en
    sentido antihorario.
\end{lemma}

\begin{proof}
    Para simplificar la notación, escribiremos que

    \begin{equation*}
        I(R) = \int_{\gamma}^{}f(z)\odif{z}.
    \end{equation*}

    Y queremos demostrar que \(I(R) = 0\) para lo cual primero necesitamos
    mostrar que \(\forall
    \varepsilon > 0,\medspace \abs{I(R)} < \varepsilon\). Además, denotaremos la
    diagonal de \(R\) por

    \begin{equation*}
        d(R) = \abs{(a + ic) - (b + id)}
    \end{equation*}

    y el perímetro de \(R\) por \(l(R)\).

    El primer paso de esta prueba es subdividir el rectángulo en cuatro subrectángulos
    iguales \(R^{(1)}, R^{(2)}, R^{(3)}\) y \(R^{(4)}\) tales que

    \begin{equation*}
        R = R^{(1)} \cup R^{(2)} \cup R^{(3)} \cup R^{(4)}\quad \wedge \quad
        \mathrm{int}(R_{i}) \cap \mathrm{int}(R_{j}) = \phi,\quad \forall i \neq j.
    \end{equation*}

    Si consideramos \(\gamma_{1}, \gamma_{2}, \gamma_{3}, \gamma_{4}\)
    parametrizaciones de
    \(\partial R^{(1)},\partial R^{(2)}, \partial R^{(3)}, \partial R^{(4)}\) en sentido
    antihorario, tenemos que

    \begin{equation*}
        I(R) = I(\partial R^{(1)}) + I(\partial R^{(2)}) +I(\partial R^{(3)})
        +I(\partial
        R^{(4)}).
    \end{equation*}

    Entre las 4 integrales, hay una que tiene el módulo más grande, sin pérdida de
    generalidad, \(\abs{I(R^{(1)})}\) es el más grande. Por la desigualdad del triángulo

    \begin{equation*}
        \abs{I(R)} \leq \sum_{i = 1}^{4} \abs{I(R^{(i)})} \leq 4 \abs{I(R^{(1)})}.
    \end{equation*}

    Por lo que \(\abs{I(R^{(1)})} \geq \tfrac{\abs{I(R)}}{4}\).

    Si repetimos el mismo proceso de subdivisión, ahora con el subrectángulo \(R_{1} =
    R^{(1)}\), tenemos que

    \begin{equation*}
        \abs{I(R_{1})} \geq \dfrac{\abs{I(R)}}{4}.
    \end{equation*}

    Siguiendo este procedimiento de manera inductiva obtenemos una sucesión
    \(\{R_{k}\}\) de
    rectángulos compactos con las siguientes propiedades:

    \begin{enumerate}[label = \arabic*.]
        \item \(R_{k + 1} \subseteq R_{k} \subseteq R\quad \forall k\in \mathbb{N}\).
        \item \(\abs{I(R_{k + 1})} \geq \abs{I(R_{k})}/4\quad \forall k\in
            \mathbb{N}\).
        \item  \(d(R_{k + 1}) = d(R_{k})/2\) y \(l(R_{k + 1}) = l(R_{k})/2,\quad
            \forall k\in \mathbb{N}\).
    \end{enumerate}

    De los puntos 2. y 3. se sigue que

    \begin{align*}
        \abs{I(R_{k})} &\leq 4 \abs{I(R_{k + 1})},\\
        d(R_{k}) &= \dfrac{d(R_{k - 1})}{2} = \dfrac{d(R_{k - 2})}{2^{2}} =
        \dfrac{d(R)}{2^{k}},\\
        l(R_{k}) &= \dfrac{l(R)}{2^{k}}.
    \end{align*}

    Entonces

    \begin{equation*}
        \abs{I(R)} \leq 4^{1}\abs{I(R_{1})} \leq 4^{2}\abs{I(R_{2})} \leq \cdots
        4^{k}\abs{I(R_{k})}\quad \forall k\in \mathbb{N}.
    \end{equation*}

    Y de la primera identidad se deduce que el diámetro de los rectángulos
    tiende a 0, i.e.

    \begin{equation*}
        \lim_{k\to \infty}d(R_{k}) = \lim_{k\to \infty}\dfrac{d(R)}{2^{k}} = 0.
    \end{equation*}

    De modo que por el teorema de los rectángulos anidados, \(\exists
    z_{0}\in \mathbb{R}\)
    tal que

    \begin{equation*}
        \bigcap_{k = 1}^{\infty} R_{k} = \{z_{0}\}.
    \end{equation*}

    Como \(z_{0}\in R \subseteq \Omega\) y \(f\) es derivable en \(z_{0}\), dado
    \(\varepsilon > 0\) arbitrario, \(\exists \delta > 0\) tal que

    \begin{align*}
        \abs`\Bigg{\dfrac{f(z) - f(z_{0})}{z - z_{0}} - f^{\prime}(z_{0})} &<
        \dfrac{\varepsilon}{d(R)l(R)},\\
        \abs{f(z) - f(z_{0}) - f^{\prime}(z_{0})\bigl(z - z_{0}\bigr)} &<
        \dfrac{\varepsilon}{d(R)l(R)}\abs{z - z_{0}},
    \end{align*}

    \(\forall z\in B_{\delta}(z_{0}) \subseteq \Omega\).

    Ahora notemos que definimos una función auxiliar que tenga primitiva, en particular

    \begin{equation*}
        g(z) = f(z_{0})z + \dfrac{1}{2}f^{\prime}(z_{0})(z - z_{0}),
    \end{equation*}

    entonces

    \begin{equation*}
        g^{\prime}(z) = f(z_{0}) + f^{\prime}(z_{0})(z - z_{0}),\quad \forall
        z\in \mathbb{C}.
    \end{equation*}

    Si suponemos que \(\tilde{\gamma}_{k}\) una parametrización de \(\partial R_{k}\) en
    sentido antihorario, por la \zcref{prop:equiv-affirmations},

    \begin{equation*}
        \int_{\tilde{\gamma}_{k}}^{}\Bigl(f(z_{0}) + f^{\prime}(z_{0})(z -
        z_{0})\Bigr)\odif{z} = 0\quad \forall k\in \mathbb{N}.
    \end{equation*}

    Además, como \(z_{0}\in R_{k}\medspace \forall k\in \mathbb{N}\) y \(d(R_{k})
    \xlongrightarrow[k\to \infty]{} 0\), existe \( N \subseteq \mathbb{N}\) tal que si
    \( k \geq N\), entonces \(R_{k} \subseteq B_{\delta}(z_{0})\).

    Con toda la información anterior, si \(k \geq N\), entonces \(\abs{z - z_{0}} \leq
        d(R_{k})\quad \forall z\in \partial R_{k} \subseteq R_{k} \subseteq
    B_{\delta}(z_{0})\),
    por lo tanto

    \begin{align*}
        \abs{I(R_{k})} &= \abs`\Bigg{\int_{\tilde{\gamma}_{k}}^{}f(z)\odif{z}},\\
        &= \abs`\Bigg{\int_{\tilde{\gamma}_{k}}^{}f(z)\odif{z} -
            \int_{\tilde{\gamma}_{k}}^{}\bigl(f(z_{0}) + f^{\prime}(z_{0})(z
        - z_{0})\bigr)\odif{z}},\\
        &= \abs`\Bigg{\int_{\tilde{\gamma}_{k}}^{}\bigl(f(z) - f(z_{0}) -
                f^{\prime}(z_{0})(z
        - z_{0})\bigr)\odif{z}},\\
        &\leq \int_{\tilde{\gamma}_{k}}^{}\abs`\Big{f(z) - f(z_{0}) -
        f^{\prime}(z_{0})\odif{z}},\\
        &\leq \int_{\tilde{\gamma}_{k}}^{}\abs`\Big{f(z) - f(z_{0}) -
        f^{\prime}(z_{0})}\abs{\odif{z}},\\
        &\leq \int_{\tilde{\gamma}_{k}}^{}\dfrac{\varepsilon}{d(R)l(R)}\abs{z -
        z_{0}}\abs{\odif{z}},\\
        &\leq \dfrac{\varepsilon
        d(R_{k})}{d(R)l(R)}\int_{\tilde{\gamma}_{k}}^{}\abs{\odif{z}},\\
        &= \dfrac{\varepsilon}{d(R)l(R)}d(R_{k})l(R_{k}).
    \end{align*}

    Usando las desigualdades e identidades obtenidas de los puntos 2. y 3.,
    concluimos que

    \begin{align*}
        \abs{I(R)} \leq 4^{k}\abs{I(R_{k})} &\leq
        4^{k}\dfrac{\varepsilon}{d(R)l(R)}d(R_{k})l(R_{k}),\\
        &= 4^{k}\dfrac{\varepsilon}{d(R)l(R)}\dfrac{d(R)}{2^{k}}\dfrac{l(R)}{2^{k}},\\
        &= \varepsilon.
    \end{align*}

    Y por lo tanto,

    \begin{equation*}
        I(R) = 0.
    \end{equation*}
\end{proof}
\end{document}
